\documentclass[14pt,a4paper,oneside]{extbook}
\usepackage[utf8]{vietnam}
\usepackage{amsmath,amssymb,mathrsfs,fancyhdr,enumerate,multirow,makecell}
\usepackage{tkz-euclide,tikz-3dplot,tikz,tkz-tab}
\usetikzlibrary{arrows.meta,calc,intersections,angles,quotes,backgrounds,shapes.geometric,patterns,through,shadows,positioning}
\usepackage{yhmath}
\usepackage[top=0.7cm, bottom=0.7cm, left=1.0cm, right=1.0cm, headheight=16pt, headsep=3mm, footskip=5mm]{geometry}
\usepackage[hidelinks,unicode]{hyperref}
\usepackage{qrcode}
\usepackage[many]{tcolorbox}
\usepackage{fontawesome}

\pagestyle{fancy}
\fancyhf{}
\fancyhead[L]{\footnotesize\bfseries\color{teal!85!black}\faBook~SỔ TAY KIẾN THỨC TOÁN 4}
\fancyhead[R]{\footnotesize\bfseries\color{blue!85!black}Lớp Toán Thầy Lĩnh -- loptoanthaylinh.vn}
\fancyfoot[C]{\footnotesize\bfseries\color{gray!80!black}-- Sổ tay kiến thức: Bài 09. Luyện tập chung về góc --}
\renewcommand{\headrulewidth}{0.4pt}

\begin{document}

% --- BANNER TIÊU ĐỀ BÀI HỌC ---
\noindent
\begin{tcolorbox}[
  enhanced,
  colback=teal!8!white,
  colframe=teal!80!black,
  arc=3mm, boxrule=1pt,
  left=3mm, right=3mm, top=2mm, bottom=2mm,
  halign=center
]
  {\Large\bfseries\color{teal!90!black}\faBookmark\ BÀI 09: LUYỆN TẬP CHUNG VỀ GÓC}\\[1.5pt]
  {\footnotesize\itshape\color{black!75} Kĩ năng đếm góc trong hình phẳng và xác định góc tạo bởi kim giờ, kim phút đồng hồ}
\end{tcolorbox}

\vspace{2.5mm}

% --- KHỐI 1 & 2: LÝ THUYẾT & KĨ NĂNG ---
\noindent
\begin{minipage}[t]{0.49\textwidth}
  \begin{tcolorbox}[
    enhanced,
    title={\footnotesize\bfseries\color{white}\faClockO\ 1. Góc trên mặt đồng hồ},
    colback=blue!2!white,
    colframe=blue!75!black,
    arc=2.5mm, boxrule=0.8pt, top=1.8mm, bottom=1.8mm, left=2mm, right=2mm
  ]
    \footnotesize
    $\bullet$ Toàn bộ vòng đồng hồ tròn là $360^\circ$ chia đều 12 khoảng số.\\
    $\rightarrow$ \textbf{Mỗi khoảng giờ tương ứng:} $360^\circ : 12 = \textbf{30}^\circ$.\\
    $\bullet$ \textbf{Lúc 3 giờ \& 9 giờ:} Cách nhau 3 khoảng $\rightarrow 30^\circ \times 3 = \textbf{90}^\circ$ (\textbf{Góc vuông}).\\
    $\bullet$ \textbf{Lúc 6 giờ:} Cách 6 khoảng $\rightarrow 30^\circ \times 6 = \textbf{180}^\circ$ (\textbf{Góc bẹt}).\\
    $\bullet$ \textbf{Lúc 1h, 2h:} Góc nhọn; \textbf{Lúc 4h, 5h:} Góc tù.
  \end{tcolorbox}
\end{minipage}%
\hfill
\begin{minipage}[t]{0.49\textwidth}
  \begin{tcolorbox}[
    enhanced,
    title={\footnotesize\bfseries\color{white}\faObjectGroup\ 2. Kĩ thuật đếm góc trong hình},
    colback=teal!2!white,
    colframe=teal!75!black,
    arc=2.5mm, boxrule=0.8pt, top=1.8mm, bottom=1.8mm, left=2mm, right=2mm
  ]
    \footnotesize
    $\bullet$ \textbf{Đếm góc tại từng đỉnh:} Xét từng đỉnh của hình đa giác, quan sát các cặp cạnh xuất phát từ đỉnh đó.\\
    $\bullet$ \textbf{Quy tắc đếm góc đơn và góc ghép:}\\
    Tại 1 đỉnh có $n$ tia chung gốc thì số góc tạo thành là: $\dfrac{n \times (n - 1)}{2}$.\\
    $\bullet$ Dùng ê-ke rà soát lần lượt từng góc để phân loại nhọn, vuông, tù.
  \end{tcolorbox}
\end{minipage}

\vspace{2.5mm}

% --- KHỐI 3: SƠ ĐỒ TƯ DUY ---
\noindent
\begin{tcolorbox}[
  enhanced,
  title={\footnotesize\bfseries\color{white}\faSitemap\ Sơ đồ tư duy: Phân loại góc tạo bởi kim giờ và kim phút đúng giờ},
  colback=yellow!2!white,
  colframe=teal!85!black,
  arc=3mm, boxrule=0.9pt, top=1.8mm, bottom=1.8mm, left=1.5mm, right=1.5mm,
  halign=center
]
  \begin{tikzpicture}[
    scale=0.92,
    every node/.style={transform shape},
    root/.style={rectangle, rounded corners=5pt, fill=teal!85!black, text=white, font=\bfseries\footnotesize, inner sep=5pt, align=center},
    b1/.style={rectangle, rounded corners=4pt, fill=blue!12, draw=blue!75!black, text=blue!90!black, font=\bfseries\scriptsize, inner sep=3.5pt, align=center},
    b2/.style={rectangle, rounded corners=4pt, fill=teal!15, draw=teal!80!black, text=teal!95!black, font=\bfseries\scriptsize, inner sep=3.5pt, align=center},
    b3/.style={rectangle, rounded corners=4pt, fill=orange!12, draw=orange!75!black, text=orange!90!black, font=\bfseries\scriptsize, inner sep=3.5pt, align=center},
    b4/.style={rectangle, rounded corners=4pt, fill=purple!12, draw=purple!75!black, text=purple!90!black, font=\bfseries\scriptsize, inner sep=3.5pt, align=center}
  ]
    \node[root] (core) at (0, 0) {\faClockO\ GÓC TRÊN MẶT ĐỒNG HỒ};

    \node[b1] (n1) at (-4.8, 0.75) {\textbf{GÓC NHỌN}\\[1pt]Lúc 1 giờ ($30^\circ$), 2 giờ ($60^\circ$)};
    \node[b2] (n2) at (4.8, 0.75) {\textbf{GÓC VUÔNG}\\[1pt]Lúc 3 giờ ($90^\circ$), 9 giờ ($90^\circ$)};
    \node[b3] (n3) at (-4.8, -0.75) {\textbf{GÓC TÙ}\\[1pt]Lúc 4 giờ ($120^\circ$), 5 giờ ($150^\circ$)};
    \node[b4] (n4) at (4.8, -0.75) {\textbf{GÓC BẸT}\\[1pt]Lúc 6 giờ ($180^\circ$)};

    \draw[teal!85!black, thick, -{Stealth[length=2mm]}] (core.north) |- (n1.east);
    \draw[teal!85!black, thick, -{Stealth[length=2mm]}] (core.north) |- (n2.west);
    \draw[teal!85!black, thick, -{Stealth[length=2mm]}] (core.south) |- (n3.east);
    \draw[teal!85!black, thick, -{Stealth[length=2mm]}] (core.south) |- (n4.west);
  \end{tikzpicture}
\end{tcolorbox}

\vspace{2.5mm}

% --- KHỐI 4: MẸO & BẪY ---
\noindent
\begin{minipage}[t]{0.54\textwidth}
  \begin{tcolorbox}[
    enhanced,
    title={\footnotesize\bfseries\color{white}\faKey\ Quy tắc đánh dấu góc khi đếm hình},
    colback=green!3!white,
    colframe=green!60!black,
    arc=2.5mm, boxrule=0.8pt, top=1.8mm, bottom=1.8mm, left=2.2mm, right=2.2mm
  ]
    \footnotesize
    $\bullet$ Đánh dấu hình cung nhỏ cho góc nhọn, hình vuông nhỏ cho góc vuông.\\
    $\bullet$ Đếm góc đơn trước, sau đó ghép 2 góc đơn liền kề xem có tạo thành góc mới không.\\
    $\bullet$ Tổng 3 góc trong một tam giác luôn bằng $180^\circ$.
  \end{tcolorbox}
\end{minipage}%
\hfill
\begin{minipage}[t]{0.44\textwidth}
  \begin{tcolorbox}[
    enhanced,
    title={\footnotesize\bfseries\color{white}\faExclamationTriangle\ Mẹo tránh bẫy trắc nghiệm},
    colback=red!2!white,
    colframe=red!70!black,
    arc=2.5mm, boxrule=0.8pt, top=1.8mm, bottom=1.8mm, left=2.2mm, right=2.2mm
  ]
    \footnotesize
    $\bullet$ \textbf{Bẫy đếm sót góc bẹt:}\\
    Hai cạnh tạo thành đường thẳng chính là một góc bẹt ($180^\circ$), học sinh rất hay bỏ sót góc này.\\
    $\bullet$ \textbf{Bẫy đồng hồ lúc 6 giờ 30 phút:}\\
    Lúc 6h30 kim giờ đã lệch khỏi số 6 $\rightarrow$ Không phải góc bẹt hay 0 độ!
  \end{tcolorbox}
\end{minipage}

\end{document}
